\chapter{Eating Disorders}

\section*{Definition and Overview}
Eating disorders are serious mental health conditions characterised by disturbances in eating behaviour, body image, and weight control, along with significant distress and impairment. The main types are anorexia nervosa, in which people severely restrict food intake; bulimia nervosa, in which episodes of binge eating are followed by purging; and binge eating disorder, in which people binge eat without purging. Eating disorders affect people of all ages, genders, and backgrounds, and they have the highest mortality of any mental illness. They are treatable, and early intervention greatly improves outcomes.

\section*{Epidemiology}
Eating disorders affect a significant number of Australians, with the most common being binge eating disorder, followed by bulimia nervosa and anorexia nervosa. They most often begin in adolescence and young adulthood but can start at any age, and they affect males as well as females, who are more commonly diagnosed. Eating disorders frequently go unrecognised, and a large proportion of people with these conditions never seek treatment, despite the availability of effective care.

\section*{Aetiology and Pathophysiology}
Eating disorders arise from a complex interaction of genetic, psychological, social, and cultural factors.
\begin{itemize}
    \item \textbf{Genetics:} there is a strong heritable component, and family history increases risk
    \item \textbf{Psychological factors:} perfectionism, low self-esteem, anxiety, and difficulties regulating emotions
    \item \textbf{Sociocultural factors:} pressure to be thin, idealised body images in media, and weight stigma
    \item \textbf{Life events:} trauma, bullying, and stressful transitions can trigger onset
    \item \textbf{Dieting:} restrictive dieting is a common gateway into disordered eating
    \item \textbf{Brain and behaviour:} disordered eating disrupts appetite regulation, and starvation itself causes physical and psychological effects that maintain the illness
    \item \textbf{Physical effects:} malnutrition affects every body system, including the heart, bones, brain, and hormones
\end{itemize}

\section*{Clinical Features}
The features vary by type of eating disorder but overlap considerably.
\begin{itemize}
    \item \textbf{Anorexia nervosa:} severe restriction of food, intense fear of gaining weight, distorted body image, and significant weight loss; often accompanied by excessive exercise
    \item \textbf{Bulimia nervosa:} recurrent episodes of binge eating followed by compensatory behaviours such as self-induced vomiting, laxative misuse, or fasting
    \item \textbf{Binge eating disorder:} recurrent binge eating without compensatory behaviours, often with feelings of shame and loss of control
    \item Preoccupation with food, weight, and body shape
    \item Avoidance of eating with others and secretive eating
    \item Physical signs: weight changes, fatigue, dizziness, digestive problems, and in women, loss of periods
    \item Mood changes, irritability, social withdrawal, and reduced concentration
    \item Wearing baggy clothes to hide weight loss or shape
\end{itemize}

\section*{Differential Diagnosis}
Other conditions can be confused with eating disorders.
\begin{itemize}
    \item Other mental health conditions, including depression, anxiety, and obsessive-compulsive disorder
    \item Avoidant/restrictive food intake disorder, in which eating is restricted without body image concerns
    \item Medical conditions causing weight loss, including thyroid disease, diabetes, gastrointestinal disease, and cancer
    \item Body dysmorphic disorder, focused on perceived appearance defects
    \item Substance use affecting appetite and weight
\end{itemize}

\section*{When to Seek Medical Care}
Early help-seeking is critical. Family and friends who suspect an eating disorder should encourage a GP review, and GPs should assess eating, weight, and psychological wellbeing routinely. Anyone with disordered eating should be assessed even if symptoms seem mild.

\textbf{Call 000 (or local emergency services) for:}
\begin{itemize}
    \item Fainting, collapse, or severe weakness
    \item Chest pain, palpitations, or irregular heartbeat
    \item Signs of severe dehydration
    \item Self-harm or thoughts of suicide
    \item Sudden severe abdominal pain, which may indicate a complication of purging
\end{itemize}

\section*{Diagnosis and Investigations}
Diagnosis is made on clinical assessment, with investigations to assess physical health.
\begin{itemize}
    \item \textbf{History:} eating behaviours, weight history, body image concerns, and mental state
    \item \textbf{Physical examination:} weight, height, body mass index, blood pressure, heart rate, and temperature
    \item \textbf{Blood tests:} electrolytes, kidney and liver function, blood counts, and blood sugar
    \item \textbf{Bone density scanning:} for people with prolonged anorexia, to assess osteoporosis risk
    \item \textbf{ECG:} to assess the heart, which can be affected by starvation and purging
    \item \textbf{Screening questionnaires:} tools that identify disordered eating patterns
    \item \textbf{Mental health assessment:} for coexisting depression, anxiety, and risk of self-harm
\end{itemize}

\section*{Management}
Treatment is multidisciplinary and addresses both the physical and psychological aspects.
\begin{itemize}
    \item \textbf{Family-based treatment:} the leading approach for adolescents with anorexia, involving parents in restoring weight
    \item \textbf{Psychological therapy:} cognitive behavioural therapy is effective for bulimia and binge eating disorder; specialist therapies address the underlying thoughts and behaviours
    \item \textbf{Nutritional rehabilitation:} dietitian input to restore weight and normalise eating
    \item \textbf{Medical monitoring:} regular review of weight, vital signs, and blood tests
    \item \textbf{Medication:} antidepressants, particularly for bulimia and binge eating disorder
    \item \textbf{Hospital care:} for severe malnutrition, medical instability, or risk of self-harm
    \item \textbf{Long-term support:} ongoing therapy and monitoring, because recovery takes time and relapse is common
\end{itemize}

\section*{Complications}
\begin{itemize}
    \item Severe malnutrition and its effects on the heart, bones, and brain
    \item Electrolyte disturbances, which can cause dangerous heart rhythms
    \item Dental erosion from vomiting
    \item Osteoporosis from prolonged low weight and hormonal disturbance
    \item Infertility and hormonal problems
    \item Depression, anxiety, and suicide risk
    \item Effects on growth and development in young people
\end{itemize}

\section*{Prognosis and Outcomes}
Eating disorders are treatable, and recovery is achievable. Outcomes are best with early intervention and family involvement, and many people make a full recovery, although the process often takes years and relapses occur. Anorexia nervosa carries the highest mortality of any mental illness, largely from medical complications and suicide, which underscores the importance of early and persistent treatment. Binge eating disorder and bulimia respond particularly well to psychological therapy.

\section*{Prevention and Self-Care}
\begin{itemize}
    \item Promote positive body image and avoid commenting on weight, in families and communities
    \item Encourage regular, flexible family meals
    \item Be alert to signs of disordered eating and seek help early
    \item Avoid restrictive dieting, especially in young people
    \item Seek treatment for anxiety, perfectionism, and low self-esteem early
    \item Challenge weight stigma and unrealistic body ideals
    \item Support loved ones with empathy rather than judgement
\end{itemize}

\section*{Sources and Further Reading}
The following reputable sources provide further information for healthcare students and patients seeking reliable, current advice about eating disorders.
\begin{itemize}
    \item Butterfly Foundation. \textit{Eating disorder support and information.} https://butterfly.org.au/
    \item National Eating Disorders Collaboration. \textit{Resources on eating disorders.} https://nedc.com.au/
    \item NHS. \textit{Eating disorders: overview and treatment.} https://www.nhs.uk/mental-health/conditions/eating-disorders/
    \item Mayo Clinic. \textit{Eating disorders: symptoms and causes.} https://www.mayoclinic.org/diseases-conditions/eating-disorders/
    \item Head to Health. \textit{Eating disorders resources.} https://www.headtohealth.gov.au/
    \item MedlinePlus. \textit{Eating disorders.} https://medlineplus.gov/eatingdisorders.html
\end{itemize}
